\documentclass[10pt,a4paper]{amsart}
\usepackage{amsmath,amssymb,amsthm,graphicx}
% For commutative diagrams we use Xy-pic
\usepackage[all]{xy}

% We use 2cell for 2-commutative diagrams.
\xyoption{2cell}
% Theorem environments.
%
\theoremstyle{plain}
\newtheorem{theorem}[subsection]{Theorem}
\newtheorem{proposition}[subsection]{Proposition}
\newtheorem{lemma}[subsection]{Lemma}

\theoremstyle{definition}
\newtheorem{definition}[subsection]{Definition}
\newtheorem{example}[subsection]{Example}
\newtheorem{exercise}[subsection]{Exercise}
\newtheorem{situation}[subsection]{Situation}

\theoremstyle{remark}
\newtheorem{remark}[subsection]{Remark}
\newtheorem{remarks}[subsection]{Remarks}

\numberwithin{equation}{subsection}

% Macros
%
\def\lim{\mathop{\mathrm{lim}}\nolimits}
\def\colim{\mathop{\mathrm{colim}}\nolimits}
\def\Spec{\mathop{\mathrm{Spec}}}
\def\Hom{\mathop{\mathrm{Hom}}\nolimits}
\def\Ext{\mathop{\mathrm{Ext}}\nolimits}
\def\SheafHom{\mathop{\mathcal{H}\!\mathit{om}}\nolimits}
\def\SheafExt{\mathop{\mathcal{E}\!\mathit{xt}}\nolimits}
\def\Sch{\mathit{Sch}}
\def\Mor{\mathop{\mathrm{Mor}}\nolimits}
\def\Ob{\mathop{\mathrm{Ob}}\nolimits}
\def\Sh{\mathop{\mathit{Sh}}\nolimits}
\def\NL{\mathop{N\!L}\nolimits}
\def\CH{\mathop{\mathrm{CH}}\nolimits}
\def\proetale{{pro\text{-}\acute{e}tale}}
\def\etale{{\acute{e}tale}}
\def\QCoh{\mathit{QCoh}}
\def\Ker{\mathop{\mathrm{Ker}}}
\def\Im{\mathop{\mathrm{Im}}}
\def\Coker{\mathop{\mathrm{Coker}}}
\def\Coim{\mathop{\mathrm{Coim}}}

% Boxtimes
%
\DeclareMathSymbol{\boxtimes}{\mathbin}{AMSa}{"02}

%
% Macros for moduli stacks/spaces
%
\def\QCohstack{\mathcal{QC}\!\mathit{oh}}
\def\Cohstack{\mathcal{C}\!\mathit{oh}}
\def\Spacesstack{\mathcal{S}\!\mathit{paces}}
\def\Quotfunctor{\mathrm{Quot}}
\def\Hilbfunctor{\mathrm{Hilb}}
\def\Curvesstack{\mathcal{C}\!\mathit{urves}}
\def\Polarizedstack{\mathcal{P}\!\mathit{olarized}}
\def\Complexesstack{\mathcal{C}\!\mathit{omplexes}}
% \Pic is the operator that assigns to X its picard group, usage \Pic(X)
% \Picardstack_{X/B} denotes the Picard stack of X over B
% \Picardfunctor_{X/B} denotes the Picard functor of X over B
\def\Pic{\mathop{\mathrm{Pic}}\nolimits}
\def\Picardstack{\mathcal{P}\!\mathit{ic}}
\def\Picardfunctor{\mathrm{Pic}}
\def\Deformationcategory{\mathcal{D}\!\mathit{ef}}
\usepackage{cleveref}
\crefname{originalreference}{original reference}{original references}
\Crefname{originalreference}{Original reference}{Original references}
\makeatletter
\newlabel{definition-locally-acyclic}{{0GJP}{source}{}{sourcepointer-1}{}}
\newlabel{definition-locally-acyclic@cref}{{[originalreference][0][]0GJP}{[1][0][]source}{}{}{}}
\makeatother
\crefname{originalreference}{original reference}{original references}
\Crefname{originalreference}{Original reference}{Original references}
\makeatletter
\newlabel{lemma-colimit}{{00D5}{source}{}{sourcepointer-2}{}}
\newlabel{lemma-colimit@cref}{{[originalreference][0][]00D5}{[1][0][]source}{}{}{}}
\makeatother
\title{A locally acyclic morphism is locally acyclic relative to every locally constant torsion sheaf of prime-to-characteristic order}
\author{The Stacks Project Authors}
\date{stacks-project:etale-cohomology.tex:local-snapshot-20261011; GNU FDL 1.2 or later}
\begin{document}
\maketitle
\noindent\textit{Editorial note.} All delivered mathematical text is the author's own, copied from the pinned Stacks Project chapter file etale-cohomology.tex; no statement or proof step was rewritten, reordered or removed.\par\medskip
\noindent\textit{Editorial note.} The counted proof is the author's complete argument: the proof is written out in full: the trivialisation at a geometric point, the filtered-colimit decomposition of the stalk into finite prime-to-characteristic groups and Z/nZ summands, and the computation of the cohomology of the affine Milnor fibres with the colimit lemma.\par\medskip
\noindent\textit{Editorial note.} The label lemma-locally-acyclic-locally-constant names the counted statement, and the proof environment immediately following it in the source is delivered verbatim as its proof.\par\medskip
\noindent\textit{Editorial note.} The counted statement is carried by the author's own lemma environment in the chapter source.\par\medskip
\noindent\textit{Editorial note.} The statement and proof are copied from the pinned chapter file; the two cited inputs are kept as external pointers with their canonical Stacks tags.\par\medskip
\noindent\textit{Editorial note.} The preamble is the chapter's own native preamble excerpt already used for this source: the Xy-pic setup (native preamble lines 14-18) and the native theorem declarations with the mathematical macros (native preamble lines 154-217).\par\medskip
\noindent\textit{Allowed standard prerequisites (editorial).}\begin{itemize}
\item Local acyclicity of f and its description through the fibres F\_\{x,t\}. Same chapter, the definition of local acyclicity and the discussion following it, giving the vanishing of the cohomology of the constant sheaves Z/nZ on the affine Milnor fibres. Supplies the hypothesis on f and the vanishing computation used for each Z/n\_\{i,j\}Z.
\item Constant sheaves and filtered colimits of cohomology. Same chapter, the lemma computing the cohomology of a constant sheaf with coefficients in a filtered colimit as the colimit of the cohomologies. Supplies the identity H\textasciicircum{}q(F\_\{x,t\}, underline M) = colim H\textasciicircum{}q(F\_\{x,t\}, underline M\_i) for the finite subgroups of M.
\end{itemize}\medskip
\ifdefined\hypertarget\hypertarget{sourcepointer-1}{}\fi
\noindent\textit{Source reference.} Original reference 0GJP: sources/stacks-etale-cohomology/pinned-original-source.tex, source lines 19110--19110. Prerequisite pointer invoked at proof line(s) 19274; it names a standard fact of the same chapter that the counted argument uses, and every step of the counted proof built on it is retained verbatim in the delivered body.\par\medskip
\ifdefined\hypertarget\hypertarget{sourcepointer-2}{}\fi
\noindent\textit{Source reference.} Original reference 00D5: sources/stacks-etale-cohomology/pinned-original-source.tex, source lines 7427--7427. Prerequisite pointer invoked at proof line(s) 19259; it names a standard fact of the same chapter that the counted argument uses, and every step of the counted proof built on it is retained verbatim in the delivered body.\par\medskip
\begin{lemma}
\label{lemma-locally-acyclic-locally-constant}
Let $f : X \to S$ be a morphism of schemes. Let $\mathcal{F}$ be a
locally constant abelian sheaf on $X_\etale$ such that for every geometric
point $\overline{x}$ of $X$ the abelian group
$\mathcal{F}_{\overline{x}}$ is a torsion group all of whose elements have
order prime to the characteristic of the
residue field of $\overline{x}$. If $f$ is locally acyclic, then $f$ is
locally acyclic relative to $\mathcal{F}$.
\end{lemma}

\begin{proof}
Namely, let $\overline{x}$ be a geometric point of $X$.
Since $\mathcal{F}$ is locally constant we see that
the restriction of $\mathcal{F}$ to $\Spec(\mathcal{O}^{sh}_{X, \overline{x}})$
is isomorphic to the constant sheaf $\underline{M}$ with
$M = \mathcal{F}_{\overline{x}}$. By assumption we can write
$M = \colim M_i$ as a filtered colimit of finite abelian groups
$M_i$ of order prime to the characteristic of the residue field of
$\overline{x}$. Consider a geometric point $\overline{t}$ of
$\Spec(\mathcal{O}^{sh}_{S, f(\overline{x})})$.
Since $F_{\overline{x}, \overline{t}}$ is affine, we have
$$
H^q(F_{\overline{x}, \overline{t}}, \underline{M}) =
\colim H^q(F_{\overline{x}, \overline{t}}, \underline{M_i})
$$
by Lemma \ref{lemma-colimit}.
For each $i$ we can write $M_i = \bigoplus \mathbf{Z}/n_{i, j}\mathbf{Z}$
as a finite direct sum for some integers $n_{i, j}$ prime to the
characteristic of the residue field of $\overline{x}$.
Since $f$ is locally acyclic we see that
$$
H^q(F_{\overline{x}, \overline{t}},
\underline{\mathbf{Z}/n_{i, j}\mathbf{Z}}) =
\left\{
\begin{matrix}
\mathbf{Z}/n_{i, j}\mathbf{Z} & \text{if} & q = 0 \\
0 & \text{if} & q \not = 0
\end{matrix}
\right.
$$
See discussion following Definition \ref{definition-locally-acyclic}.
Taking the direct sums and the colimit we conclude that
$$
H^q(F_{\overline{x}, \overline{t}}, \underline{M}) =
\left\{
\begin{matrix}
M & \text{if} & q = 0 \\
0 & \text{if} & q \not = 0
\end{matrix}
\right.
$$
and we win.
\end{proof}

\end{document}
